% =====================================================================
% Revision for identifiability_rank_bound (v5): separate what is ESTABLISHED
% from what is CONJECTURED, and replace the "capacity formula" emphasis with
% the end-to-end Jacobian singular spectrum + a falsifiable experiment.
%
% HOW TO USE: drop the FRAMING box near the top (after the abstract / "read
% this first"); replace the old feasibility sections (the "intersection view"
% / "operative calculation" / "capacity" sections, i.e. old sec.9-11) with the
% sections below; put the BOTTOM-LINE box at the very end. Boxes are written
% as \begin{revbox}...\end{revbox}; either define that (a thin tcolorbox
% wrapper) or replace with your existing highlight environment.
%
% Citations are inline placeholders [Author, venue]; give them real bib keys.
% Numbers quoted from the literature are marked % VERIFY.
% =====================================================================

% ---- optional lightweight box (or reuse your own) ----
% \usepackage[most]{tcolorbox}
% \newtcolorbox{revbox}{colback=blue!4,colframe=blue!30,boxrule=0.4pt,arc=1pt,left=6pt,right=6pt,top=4pt,bottom=4pt}

% =====================================================================
% FRAMING BOX  (place near the top)
% =====================================================================
\begin{revbox}
\textbf{What this note does and does not establish.}
The rank theorem (Secs.~3--5) is exact, but only about an \emph{information-mixing
mechanism}, under the assumption that the gate matrix $G$ is known and frozen. It is
\emph{not} a reconstruction attack: a real attacker does not know $G$, and $G$ itself
depends on the unknown images. So the rank condition is a necessary structural fact,
not a proof of leakage. The constructive half of the note is therefore a
\emph{conjecture with a test}, not a theorem:
\begin{center}\itshape
This note does not prove that an ordinary LoRA adapter leaks its training images.
It gives a plausible mechanism for why it could, and --- more importantly --- a
concrete experiment that decides whether the information needed for reconstruction is
actually present in the released adapter.
\end{center}
\end{revbox}

% =====================================================================
\section{Read the rank theorem correctly: mechanism, not attack}
Sections~3--5 prove one exact thing. If the gradient is $\Omega=GX^{\top}$, then every
row of $\Omega$ is a fixed mixture $\sum_i G_{ki}\,x_i^{\top}$ of the private images, and
the images survive \emph{in separable form} iff the neurons supply enough independent
mixtures --- precisely $\rank(G)=\rank(M)=\rho$. When $\rho<N$ some images collapse into
a single shared sum and cannot be told apart. This is exact under the
\emph{frozen-known-$G$} reading.

The limitation is equally exact: this is not yet an attack. The real attacker neither
knows $G$ nor can hold it fixed, because $G=\sigma'(\langle w_k,x_i\rangle)$ is a
function of the very images being sought. So treat Secs.~3--5 as characterising the
\emph{information-mixing mechanism} of the measurement, and read $\rho\ge N$ as a
necessary condition on the first-order signal --- the floor, not the verdict.

% =====================================================================
\section{Fine-tuning as a measurement system}
Write the whole experiment as one map. A latent $z=(z_1,\dots,z_N)$ generates images
$x=(x_1,\dots,x_N)=g(z)$, which are fine-tuned into an adapter:
\[
z \;\xrightarrow{\;g\;}\; x \;\xrightarrow{\text{LoRA training}}\; (A_T,B_T)
\;=:\; F(z).
\]
$F$ is a measurement system with \emph{two} lossy stages, and both help the attacker:
\begin{itemize}\itemsep2pt
\item \textbf{Images are not arbitrary pixels.} Moving between realistic images requires
coordinated changes (pose, lighting, shape, texture), so the data occupies a manifold
$\mathcal M_X$ of local dimension $k\ll d$.
\item \textbf{Gradients of real images are not arbitrary either.} The set
$\mathcal M_G=\{\nabla_\theta L(X):X\ \text{realistic}\}$ is itself a thin subset of the
ambient gradient space.
\end{itemize}
The whole privacy question is then sharp: \emph{is $F$ locally injective and
well-conditioned on $\mathcal M_X$?} LoRA may destroy vast amounts of \emph{arbitrary}
gradient information while still separating points of $\mathcal M_G$ --- a restricted
inverse, not a full one.

% =====================================================================
\section{Dimension count $\to$ Jacobian rank $\to$ singular spectrum}
The manifold argument gives a first, weak checkpoint. If image $i$ has $k_i$ local
degrees of freedom, the private unknown has $q=\sum_i k_i$ directions, while a rank-$\rho$
measurement spans at most $d\rho$ independent directions, so a regular locally invertible
map needs
\[
d\rho \;\gtrsim\; q \;=\;\textstyle\sum_i k_i .
\]
This is worth stating plainly for what it is: a \textbf{dimensional plausibility check},
necessary and nothing more. Having enough dimensions does not mean they point in useful
directions. The map
$A=\left(\begin{smallmatrix}1&0\\0&10^{-10}\end{smallmatrix}\right)$ preserves rank $2$,
yet the second coordinate is gone for any practical purpose. So the honest progression is
\[
\boxed{\text{dimension count}}\;\longrightarrow\;
\boxed{\text{Jacobian rank}}\;\longrightarrow\;
\boxed{\text{singular spectrum}} ,
\]
and only the last object measures usable leakage. The earlier $d\rho\ge Nk$ inequality
should be read as sizing the \emph{null space}, not as a capacity law.

% =====================================================================
\section{The object that actually decides it: the end-to-end Jacobian}
Define the Jacobian of the entire data$\to$adapter map,
\[
J_{\mathrm{full}} \;=\;
\frac{\partial\,\operatorname{vec}(A_T,B_T)}{\partial(z_1,\dots,z_N)} .
\]
Its meaning is concrete and testable. Take one tiny realistic change to one training
image, retrain the LoRA, and ask whether the released adapter changed --- and whether a
\emph{different} realistic change moves it \emph{differently}. If every private-data
direction produces a distinguishable adapter change, the adapter has locally retained the
information; if some direction alters the images substantially but barely moves the
adapter, that direction has been erased. The singular values of $J_{\mathrm{full}}$ read
this off directly, and the note's strongest empirical prediction is
\[
\boxed{\;\sigma_{\min}(J_{\mathrm{full}})\ \text{collapses in the same regimes where
reconstruction collapses.}\;}
\]
This is a far cleaner claim than ``rank-$r$ LoRA allows $r$ examples'' or any
$d\rho/N$ capacity number.

% =====================================================================
\section{The experiment that decides it}
Do not begin by estimating the intrinsic dimension of Flowers --- that swaps the open
question for a second open question. Instead make $k$ \textbf{known by construction}:
take a generator $g:\mathbb R^{k}\to\text{images}$ and restrict each training image to
$k\in\{4,8,16,32,\dots\}$ latent directions, so the private degrees of freedom are exactly
$q=Nk$ (or, precisely, $\rank J_g$, which is directly measurable).

For each $N$: sample $z\to$ train LoRA $\to$ release $Y=(A_T,B_T)$; compute the singular
spectrum of $J_{\mathrm{full}}$; and \emph{separately} run the actual inversion attack on
$Y$. Plot $N$ against $\sigma_{\min}(J_{\mathrm{full}})$ against reconstruction quality
(LPIPS). Three outcomes, all informative:
\begin{description}\itemsep2pt
\item[World A (the dream).] $\sigma_{\min}$ and reconstruction collapse at the same $N$:
an identifiability transition, with the theory explaining the attack.
\item[World B.] $J_{\mathrm{full}}$ stays well-conditioned but the attack fails: the
adapter \emph{does} preserve the information; the inversion/decoder is the bottleneck.
\item[World C.] $J_{\mathrm{full}}$ collapses yet the decoder still emits plausible
images: the decoder is \emph{hallucinating} from its prior rather than recovering
information forced by the adapter --- a critical privacy distinction.
\end{description}
Distinguishing B from C is exactly what a rank/capacity number cannot do and this
experiment can.

% =====================================================================
\section{The figure: what LoRA remembers vs.\ forgets}
Once $J_{\mathrm{full}}=U\Sigma V^{\top}$ is computed, take the weakest right-singular
vector $v_{\min}$ and perturb the data, $z'=z+\epsilon\,v_{\min}$; render the altered
images; train both datasets and compare adapters. If the story holds, $v_{\min}$ is a
\emph{meaningful change to the private images that the released adapter almost cannot
see}, while $v_{\max}$ is a comparable image-space change that moves the adapter
enormously. That turns the spectrum into a picture --- ``a private-data change LoRA
remembers'' beside ``one it forgets'' --- a compelling thesis figure.

% =====================================================================
\section{Localising the activation effect}
The same language sharpens the activation study. Compute progressively richer Jacobians
along the pipeline:
\[
J_{\mathrm{grad}}=\frac{\partial\,\nabla_\theta L}{\partial z},\qquad
J_{\text{LoRA-step}},\qquad
J_{\mathrm{full}},
\]
for ReLU, GELU, SiLU, Softplus. Then ``GELU fails'' becomes a locatable claim:
\begin{itemize}\itemsep1pt
\item If $J_{\mathrm{grad}}$ is already ill-conditioned, the activation destroys data
directions at the \emph{gradient} stage.
\item If $J_{\mathrm{grad}}$ is healthy but the LoRA-stage spectrum collapses, the loss is
in \emph{how LoRA observes} the gradient.
\item If even $J_{\mathrm{full}}$ is healthy but the attack fails, the ReLU/GELU gap is an
\emph{algorithmic inversion} phenomenon, not fundamental information loss.
\end{itemize}
This unifies the activation, anchor, and gradient-bridge tracks into one experimental
language.

% =====================================================================
\section{The gradient bridge as a restricted inverse}
The measurement view explains what the bridge could be exploiting: LoRA may be
non-invertible on \emph{arbitrary} gradients yet nearly invertible on the thin set
$\mathcal M_G$ of gradients that real images actually produce. This is consistent with
R2F's premise (a decoder learns full-gradient directions from LoRA gradients via a proxy
relationship) --- but R2F targets \emph{unlearning}, not training-image recovery
[Liu et al., 2025], and a good decoder does \emph{not} prove LoRA preserved the
full-gradient coordinates: it may be predicting missing coordinates from the gradient
prior, exactly like an image model inpainting detail.

That is testable. Run the bridge, then against controls that (i) shuffle adapters between
examples, (ii) replace the LoRA observation with matched-norm noise, (iii) remove specific
LoRA modules, (iv) feed only progressively smaller pieces of the adapter. If it still
reconstructs, the public-data prior is doing the work. More decisively, measure the
observation Jacobian $\partial Y/\partial z$ on the real gradient manifold: if
\[
\rank\!\left(\tfrac{\partial Y}{\partial z}\right)\;\approx\; Nk
\quad\text{despite large null spaces in the ambient full-gradient space,}
\]
that is direct evidence for the restricted-inverse story --- ``LoRA is non-invertible on
arbitrary gradients but nearly invertible on naturally occurring ones'' --- a far stronger
statement than ``we trained a decoder that predicts gradients.''

% =====================================================================
\section{Two confounds to rule out before believing the spectrum}
\paragraph{Local vs.\ global.} $J_{\mathrm{full}}$ answers a \emph{local} question. Run
inversion twice: from $z_{\mathrm{init}}=z^\star+\text{small noise}$ (does the
neighbourhood recover? --- should track $J_{\mathrm{full}}$), and from random init (can an
attacker find the point globally?). If local succeeds while global fails, the Jacobian
theory is intact and the obstacle is optimisation, not information.
\paragraph{Seed vs.\ signal.} A deterministic $Y=F(X;\xi)$ fixes the training randomness
$\xi$; a real attacker may not know minibatch order, augmentation, or dropout. Compare the
\emph{data signal} $\|\mathbb E[Y\mid X+\delta X]-\mathbb E[Y\mid X]\|$ against the
\emph{training noise} $\sqrt{\mathbb E\|Y-\mathbb E[Y\mid X]\|^2}$. If flipping one
expression moves the adapter less than reshuffling the minibatch does, deterministic
$J_{\mathrm{full}}$ overstates usable leakage; if the data signal dominates, the privacy
argument is much stronger. Not experiment~\#1, but essential once the deterministic story
holds.

% =====================================================================
\section{The \texorpdfstring{$k$}{k} you are allowed to use}
A subtle but important point: you do not get to plug a natural-image intrinsic-dimension
estimate into the formula. If the attack optimises a $512$-dim generator latent $z$ with
$\rank J_g(z)=512$ locally, the unknown has $512$ local directions, regardless of any
paper reporting $\sim\!40$ for ImageNet. The relevant $k$ is the dimension of the
\emph{hard prior the attacker is actually restricted to}, $k=\rank J_g$. That is another
reason to control $k$ by construction first, and only later ask what effective $k$ a real
decoder imposes.

% =====================================================================
\section{Where this sits in the literature}
Malicious-server results establish that low rank is not itself a privacy guarantee, but
answer a different question. MineGrad [AISTATS 2026] and the CVPR'25 PEFT attack recover
data from LoRA/PEFT gradients only after the server \emph{maliciously} alters the
pretrained model / adapter setup so the data is encoded into the shared gradients. They
show ``LoRA's low rank does not make inversion impossible,'' not ``an ordinary honest
final checkpoint contains enough to reconstruct its images.''

On the honest-checkpoint side, neighbouring evidence is suggestive but not recovery: final
LoRA weights reveal fine-tune \emph{dataset size} to within $\sim\!0.36$ images (DSiRe)
% VERIFY: DSiRe MAE ~0.36
and weight-space studies show customised LoRA weights carry strong signal about visual
identity and semantic attributes. These make ``the final adapter encodes substantial
information'' credible, but stop short of exact training-image recovery. That gap is
precisely the question this note is built to test.

Intrinsic-dimension estimates support only the \emph{qualitative} premise. Pope et al.\
report ImageNet intrinsic dimensions of roughly $26$--$43$ by nearest-neighbour estimators
% VERIFY
against $\sim\!10^5$ pixels; but estimators disagree sharply --- on MNIST a diffusion-based
estimator gives $\sim\!152$ while nearest-neighbour MLE gives $\sim\!13$--$14$.
% VERIFY
So ``natural images occupy far less than pixel space'' is believable; ``the relevant $k$
for Flowers is $37$'' is \emph{not} established --- which is exactly why $k$ must be
controlled, not assumed.

% =====================================================================
% BOTTOM-LINE BOX  (place at the very end)
% =====================================================================
\begin{revbox}
\textbf{What we actually know, and the next experiment.}
Fine-tuning is a measurement system: gradients mix training examples; low-dimensional data
structure means severe ambient compression need not destroy identifiability; LoRA applies
a further lossy measurement; and the question for a real released adapter is whether the
complete measurement is locally injective and well-conditioned \emph{on the set of
realistic datasets}. The singular spectrum of the end-to-end data$\to$adapter Jacobian is
a direct way to measure that. The toy rank theorem explains why this is conceivable; the
intrinsic-dimension literature explains why compression alone does not rule it out; the
malicious-server attacks show low rank is not a guarantee; final-LoRA studies show the
weights encode surprising properties of the fine-tuning set. But the central empirical
claim is still open:
\begin{center}\itshape
Does an ordinarily trained final LoRA preserve the private-image directions well enough
for actual reconstruction?
\end{center}
The highest-value next step is not more theory, and not estimating Flowers' manifold
dimension, but:
\[
\boxed{\textbf{control } k \;\to\; \textbf{sweep } N \;\to\;
\textbf{measure } \sigma(J_{\mathrm{full}}) \;\to\; \textbf{run inversion}}
\]
and test whether spectral collapse predicts reconstruction collapse. If it does, the note
becomes an experimentally falsifiable theory of LoRA leakage.
\end{revbox}
